(1) 利用$\sigma_x^2=I$，自由演化算符$U(\delta t)=\cos(\Omega\delta t/2)I-i\sin(\Omega\delta t/2)\sigma_x$。由$\langle g|\sigma_x|g\rangle=0$，单步基态存活振幅为$\cos(\Omega\delta t/2)$，概率$p=\cos^2(\Omega\delta t/2)$。单步偏离1的量级是$(\delta t)^2$，不是线性时间。
(2) 选择性条件要求每一步结果都为$g$，各步状态重置到$|g\rangle$，因此$P_{\rm no\,jump}=p^N=[\cos^2(\Omega t/(2N))]^N$。展开$\ln\cos^2x=-x^2+O(x^4)$得$\ln P_{\rm no\,jump}=-\Omega^2t^2/(4N)+O(N^{-3})$，所以固定$t$下$N\to\infty$有$P_{\rm no\,jump}\to1$，即量子Zeno冻结。
(3) 非选择性测量也允许从$g$跳到$e$再返回。每步状态满足二态转移矩阵$\mathsf P=\begin{pmatrix}p&1-p\\1-p&p\end{pmatrix}$，本征值为1及$2p-1=\cos(\Omega\delta t)$。初始基态后第$N$次测量得到$g$的总概率$P_g^{(N)}=[1+\cos^N(\Omega t/N)]/2$，它与从未跳跃的$p^N$不同，但两者在连续投影极限都趋1。
(4) 若改用持续$\sigma_z$退相干，设$\dot\rho=-i(\Omega/2)[\sigma_x,\rho]+(\Gamma/2)(\sigma_z\rho\sigma_z-\rho)$，非对角元以速率$\Gamma$衰减。Bloch变量$z=\operatorname{Tr}(\rho\sigma_z)$满足$\ddot z+\Gamma\dot z+\Omega^2z=0$，当$\Gamma\gg\Omega$时慢模衰减率为$\Omega^2/\Gamma+O(\Omega^4/\Gamma^3)$，有效激发交换被强测量抑制。连续退相干与离散理想投影的定量参数需单独对应，不能直接把$\Gamma$设成$N/t$而忽略测量模型。
